\chapter*{Resumo}

\par Este trabalho apresenta o desenvolvimento de uma Plataforma de Engenharia Self-Service com o objetivo de abstrair a complexidade na provisão e configuração de ambientes de infraestrutura para pesquisa e desenvolvimento. Tradicionalmente, a criação de máquinas virtuais (VMs) ou contêineres, bem como a configuração de softwares e serviços específicos, demanda tempo e conhecimento técnico especializado do usuário final. A solução desenvolvida, denominada ProxLab, transforma esse processo manual em um modelo sob demanda: por meio de um portal web, o usuário (aluno ou pesquisador) especifica suas necessidades de infraestrutura, e o backend orquestra a criação de máquinas virtuais no Proxmox VE por meio de Terraform (Infraestrutura como Código) ou o provisionamento de contêineres de aplicação em um cluster Kubernetes (k3s) dedicado, com quotas por perfil de usuário e acesso interativo via terminal web. Essa abordagem de Engenharia de Plataforma capacita pesquisadores que não possuem conhecimento aprofundado em infraestrutura a provisionarem seus próprios ambientes de trabalho de forma ágil, reduzindo a dependência de equipes de operação e a ocorrência de erros manuais. Os resultados obtidos demonstram a viabilidade técnica da abordagem, evidenciando ganhos de autonomia e redução do tempo de provisionamento em relação ao processo manual.

\vspace{2em}

\textbf{Palavras-chave:} Engenharia de Plataforma, Self-Service, Infraestrutura como Código (IaC), Terraform, Kubernetes, Proxmox, Automação de Infraestrutura.

\vfill

% ---------------------------------------------------------------------------- %